\documentclass[12pt, letterpaper]{article}

\usepackage[utf8]{inputenc}
\usepackage[T1]{fontenc}
\usepackage[margin=1in]{geometry}
\usepackage{setspace}
\usepackage{graphicx}
\usepackage{booktabs}
\usepackage{amsmath, amssymb}
\usepackage{hyperref}
\usepackage[style=apa, backend=biber]{biblatex}

\addbibresource{references.bib}
\onehalfspacing

\title{Family-as-CAS Calibration:\\Multi-Agent Relational Response Simulation Using FAD and SCORE-15 Benchmarks}
\author{Ryan Carr}
\date{\today}

\begin{document}

\maketitle

\begin{abstract}
TODO: Summarize the Family-as-CAS extension and its calibration against FAD and SCORE-15 family-system benchmarks.
\end{abstract}

\section{Introduction}
TODO: Explain why a family-system layer provides an intermediate calibration layer between individual and organizational assessment.

\section{Family Systems as Complex Adaptive Systems}
TODO: Review family functioning, relational dynamics, multi-informant perspectives, and longitudinal change.

\section{Simulation Architecture}
TODO: Describe the multi-agent family simulation and respondent perspective model.

\section{Calibration Benchmarks}
TODO: Define FAD multidomain and multi-informant targets plus SCORE-15 longitudinal-change targets.

\section{Results}
TODO: Present calibration outcomes, divergences, and uncertainty carried forward.

\section{Discussion}
TODO: Interpret the family-system layer as a methodological bridge rather than direct validation of the OHA.

\section{Conclusion}
TODO: State the contribution to staged calibration for brief survey instruments.

\printbibliography

\end{document}
